% ---
% id: statistics/binomial-box/thm:monotonicity-limits
% kind: thm
% title: "Monotonicity in the budget and limiting cases"
% summary: >-
%   The acceptance set shrinks as the budget grows under monotone allocation
%   and endpoint rules, so no infimum over it, h_low included, decreases.
%   Zero-budget and one-letter limits are explicit; at a fixed positive
%   allocation the exact rule's set tends to the envelope's box as n^(−1/2),
%   unnested in n.
% prominence: standard
% status: accepted
% uses: [def:endpoint-rule, def:monotone-rule, def:interval, def:construction, lem:guard, thm:completeness-bound, lem:cheb, thm:box-construction, statistics/def:envelope, lem:widening, statistics/def:conservative-evaluation]
% mentions: [prop:rule-exact, prop:rule-bernstein, prop:bracketed-endpoints, meat-core/def:feasible]
% assumes: [A-ALPHABET]
% refs: []
% provenance: newly-derived
% lean:
% ---

Parts (i) and (ii) need the endpoint rule itself to be
\tref{monotone in the allocation}, and which rules are is each rule's own
statement: the \tref[exact binomial quantile rule]{exact rule} with its
sweep evaluation (Proposition~\ref{prop:rule-exact}\,(ii), and (vii) where
a tail is decided against an upper bound) and the Bernstein
rule (Proposition~\ref{prop:rule-bernstein}\,(ii)) are, while the enclosure
evaluation of the exact rule (Proposition~\ref{prop:bracketed-endpoints})
does not claim it, so (i) and (ii) are not asserted for it. The monotonicity
of $\hlow$ in (ii) is a trade-off, not an improvement: the same change of
budget makes the honest run abort no less often. A one-letter acceptance
alphabet, the case of (vi), imposes no constraint on the \tref{consumer} and
carries no information; it is a degenerate case to be recognised, not a
regime to be used.

\begin{theorem}\label{thm:monotonicity-limits}
Fix $\Chatset$, an \tref{honest envelope} $({\honest}^{-},{\honest}^{+})$
with \tref{envelope set} $\mathsf{H}$ (Definition~\ref{def:envelope}; a single
law $\honest$ is the \tref{degenerate envelope}, $\mathsf{H}=\{\honest\}$),
$n$, $\Chatcon$ and an \tref{admissible endpoint rule} $R$
(Definition~\ref{def:endpoint-rule}) under a \tref{conservative evaluation}
(Definition~\ref{def:conservative-evaluation}), write $p^{\mp}_{\hat c}$ for the
envelope's values as there, and write $\acSet(e)$ for the output of
$\mathsf{Acc}_R$ at budget $e$ under a fixed monotone allocation rule
(Definition~\ref{def:monotone-rule}).
\begin{enumerate}[label=(\roman*),leftmargin=2.2em]
\item \emph{(budget monotonicity: the set shrinks as the budget grows)} If
$R$ is monotone in the allocation and $0\le e_1\le e_2<1$, then
$\acSet(e_2)\subseteq\acSet(e_1)$.
\item \emph{(consequence for any consumer infimum)} Under the hypothesis of
(i), for every function $\Psi:\simplex{\Chatset}\to[-\infty,+\infty]$ and
$e_1\le e_2$,
\begin{equation}\label{eq:inf-mono}
  \inf_{\mathbf{q}\in\acSet(e_1)}\Psi(\mathbf{q})
  \ \le\ \inf_{\mathbf{q}\in\acSet(e_2)}\Psi(\mathbf{q}) .
\end{equation}
In particular, taking
$\Psi(\mathbf{q})=\inf_{\omega\in\Feas}O_{\alpha}(\mathbf{q},\omega)$, the
quantity $\hlow$, which \xref{meat-core/def:feasible} defines as
$\inf_{\mathbf{q}\in\acSet}\inf_{\omega\in\Feas}O_{\alpha}$, is
\emph{non-decreasing} in $\ecomp$: a weaker completeness guarantee buys a
larger or equal entropy bound.
\item \emph{(limit $\ecomp=0$, every rule)} At $e=0$ every
$\varepsilon^{\pm}(\hat c)=0$ and
\begin{equation*}
\begin{aligned}
  \acSet(0)=\big\{\mathbf{q}\in\simplex{\Chatset}:\
  &q(\hat c)=0\ \text{for }\hat c\in\Chatcon\text{ with }p^{+}_{\hat c}=0,\\
  &q(\hat c)=1\ \text{for }\hat c\in\Chatcon\text{ with }p^{-}_{\hat
  c}=1\big\},
\end{aligned}
\end{equation*}
with $\varepsilon^{\star}=0$. If $p^{+}_{\hat c}>0$ and $p^{-}_{\hat c}<1$ for
every $\hat c\in\Chatcon$ (for the degenerate envelope: if
$0<\honest(\hat c)<1$ for every $\hat c\in\Chatcon$), then
$\acSet(0)=\simplex{\Chatset}$: a zero completeness budget buys the whole
simplex and no constraint at all.
\item \emph{(limit $n\to\infty$, the exact rule)} Let $R$ be the exact rule
of Definition~\ref{def:interval}, let $\hat c\in\Chatcon$ and write
$p^{\mp}=p^{\mp}_{\hat c}$. Then
\begin{equation}\label{eq:width}
  0\ \le\ q^{+}(\hat c)-p^{+}\ \le\
  \sqrt{\frac{p^{+}(1-p^{+})}{n\,\varepsilon^{+}(\hat c)}}+\frac1n,
  \qquad
  0\ \le\ p^{-}-q^{-}(\hat c)\ \le\
  \sqrt{\frac{p^{-}(1-p^{-})}{n\,\varepsilon^{-}(\hat c)}}+\frac1n ,
\end{equation}
the first whenever $0<p^{+}<1$ and $\varepsilon^{+}(\hat c)>0$, the second
whenever $0<p^{-}<1$ and $\varepsilon^{-}(\hat c)>0$.
Consequently, if $\Chatcon=\Chatset$, $0<p^{\mp}_{\hat c}<1$ and
$\varepsilon^{\pm}(\hat c)>0$ for all $\hat c$, and the allocation is held
fixed as $n$ grows (for instance the uniform rule at a fixed $e>0$), then
every $\mathbf{q}\in\acSet$ satisfies
$p^{-}_{\hat c}-O(n^{-1/2})\le q(\hat c)\le p^{+}_{\hat c}+O(n^{-1/2})$
uniformly in $\hat c$: $\acSet$ lies coordinatewise within $O(n^{-1/2})$ of
the box $\prod_{\hat c}[p^{-}_{\hat c},p^{+}_{\hat c}]$, and it contains
$\mathsf{H}$ by Theorem~\ref{thm:box-construction}\,(iii). For the degenerate envelope,
$\sup_{\mathbf{q}\in\acSet}\|\mathbf{q}-\honest\|_{\infty}=O(n^{-1/2})$ and
$\acSet$ converges to $\{\honest\}$ in Hausdorff distance.
\item \emph{(no monotonicity in $n$, with a counterexample; the exact rule)}
For the exact rule, $\acSet$ is \emph{not} monotone in $n$, already for a
degenerate envelope. Take $\Chatset=\{\hat c_0,\hat c_1\}$, the degenerate
envelope at $\honest=(\tfrac12,\tfrac12)$, $\Chatcon=\{\hat c_0\}$ and
$\varepsilon^{-}(\hat c_0)=\varepsilon^{+}(\hat c_0)=\tfrac25$ (so
$\ecomp=\tfrac45$). Then
\begin{equation*}
\begin{aligned}
  n&=2:\quad [k^{-},k^{+}]=[1,1],\quad
   \acSet=\{(\tfrac12,\tfrac12)\};\\
  n&=3:\quad [k^{-},k^{+}]=[1,2],\quad
   \acSet=\{\mathbf{q}\in\simplex{\Chatset}:\tfrac13\le q(\hat
   c_0)\le\tfrac23\},
\end{aligned}
\end{equation*}
so $\acSet(n{=}2)\subsetneq\acSet(n{=}3)$: the set is strictly larger at the
larger round count.
\item \emph{(degenerate alphabet $|\Chatset|=1$, every rule)} If
$\Chatset=\{\hat c_0\}$ then, for every $n\ge1$, every $\ecomp\in[0,1)$,
every $\Chatcon\subseteq\Chatset$ and every admissible allocation,
\begin{equation*}
  {\honest}^{+}(\hat c_0)=1,
  \qquad
  \mathsf{H}=\acSet=\simplex{\Chatset}=\{\delta_{\hat c_0}\},
\end{equation*}
and $\varepsilon^{\star}=0$ whenever ${\honest}^{-}(\hat c_0)=1$, in
particular for the degenerate envelope, which is then
$\honest=\delta_{\hat c_0}$. The emitted set is therefore constant in all of
the construction's parameters on a one-letter alphabet, being at once the
whole simplex and a single point; for the degenerate envelope so are its box
data, and completeness is perfect. For a non-degenerate envelope
(${\honest}^{-}(\hat c_0)<1$, which \eqref{eq:envelope} permits here) the
lower box bound $q^{-}(\hat c_0)$ and the reported $\varepsilon^{\star}$ may be
positive and may vary with the parameters; the reported figure is still a
valid bound (by Theorem~\ref{thm:completeness-bound}), merely a loose one, the true honest
failure probability being $0$.
\end{enumerate}
\end{theorem}

\begin{proof}
\emph{(i).} Fix $\hat c\in\Chatcon$, write $p^{\mp}=p^{\mp}_{\hat c}$ and
$\varepsilon^{\pm}_i\coloneqq\varepsilon^{\pm}_{(e_i)}(\hat c)$, so
$\varepsilon^{\pm}_1\le\varepsilon^{\pm}_2$ by
Definition~\ref{def:monotone-rule}. Since $R$ is monotone in the allocation
(Definition~\ref{def:endpoint-rule}),
$R^{-}(n,p^{-},\varepsilon^{-}_1)\le R^{-}(n,p^{-},\varepsilon^{-}_2)$ and
$R^{+}(n,p^{+},\varepsilon^{+}_1)\ge R^{+}(n,p^{+},\varepsilon^{+}_2)$. Both
guards in \eqref{eq:guard} are monotone in their first argument, their second
arguments $\lfloor np^{-}\rfloor$ and $\lceil np^{+}\rceil$ not depending on
the budget, whence $k^{-}_1(\hat c)\le k^{-}_2(\hat c)$ and
$k^{+}_1(\hat c)\ge k^{+}_2(\hat c)$,
that is, $[k^{-}_2,k^{+}_2]\subseteq[k^{-}_1,k^{+}_1]$. The endpoints
$k^{\mp}_1$ at the smaller budget are therefore integer endpoints of the kind
Lemma~\ref{lem:widening}\,(iv) admits for the output $\acSet(e_2)$:
$0\le k^{-}_1(\hat c)\le k^{-}_2(\hat c)$ and
$k^{+}_2(\hat c)\le k^{+}_1(\hat c)\le n$, the outer bounds by
\eqref{eq:ordering} at $e_1$ and the inner ones as just shown. The set
\eqref{eq:widened-set} they give is $\acSet(e_1)$, since both sets are built
by \eqref{eq:boxbounds} on the same $\Chatcon$, with the bounds $[0,1]$ off
$\Chatcon$ at both budgets. Lemma~\ref{lem:widening}\,(i) gives
$\acSet(e_2)\subseteq\acSet(e_1)$.

\emph{(ii).} \eqref{eq:inf-mono} is \eqref{eq:widened-inf} of
Lemma~\ref{lem:widening}\,(iii) applied to the inclusion of (i), with
$\acSet(e_2)$ in the role of $\acSet$ and $\acSet(e_1)$ in that of $\acSet'$,
and it holds
for arbitrary $\Psi$ with values in the extended reals, no continuity or
measurability required. The specialization substitutes the definition of
$\hlow$ as an infimum over $\acSet$ and changes nothing else.

\emph{(iii).} The allocation constraint at $e=0$ forces every
$\varepsilon^{\pm}(\hat c)=0$, these being non-negative with sum at most $0$.
Fix $\hat c\in\Chatcon$ and write $p^{\mp}=p^{\mp}_{\hat c}$; by
\eqref{eq:guard} each side is computed from its own envelope value alone.

\emph{Lower side}, $R^{-}=R^{-}(n,p^{-},0)$. If $p^{-}\in(0,1)$:
Definition~\ref{def:endpoint-rule}\,(b) gives $R^{-}=0$, so
$k^{-}(\hat c)=\min\{0,\lfloor np^{-}\rfloor\}=0$. If $p^{-}=0$:
Definition~\ref{def:endpoint-rule}\,(c) gives $R^{-}=0$ and
$k^{-}(\hat c)=\min\{0,0\}=0$. If $p^{-}=1$: (c) gives $R^{-}=n$ and
$k^{-}(\hat c)=\min\{n,n\}=n$. So the lower bound is $q^{-}(\hat c)=0$ unless
$p^{-}=1$, where it is $q^{-}(\hat c)=1$.

\emph{Upper side}, $R^{+}=R^{+}(n,p^{+},0)$. If $p^{+}\in(0,1)$: (b) gives
$R^{+}=n$, so $k^{+}(\hat c)=\max\{n,\lceil np^{+}\rceil\}=n$. If $p^{+}=0$:
(c) gives $R^{+}=0$ and $k^{+}(\hat c)=\max\{0,0\}=0$. If $p^{+}=1$: (c) gives
$R^{+}=n$ and $k^{+}(\hat c)=\max\{n,n\}=n$. So the upper bound is
$q^{+}(\hat c)=1$ unless $p^{+}=0$, where it is $q^{+}(\hat c)=0$.

On $\simplex{\Chatset}$ the bound $q(\hat c)\ge1$ is $q(\hat c)=1$ and the
bound $q(\hat c)\le0$ is $q(\hat c)=0$, and the bounds $[0,1]$ are vacuous;
assembling the cases in \eqref{eq:accset} gives the displayed set, and it is
the whole simplex when no $\hat c\in\Chatcon$ has $p^{+}_{\hat c}=0$ or
$p^{-}_{\hat c}=1$. Finally $\varepsilon^{\star}=0$. A lower term of
\eqref{eq:completeness} is $\Pr[\Ncount{\hat c}^{-}<0]=0$ when
$k^{-}(\hat c)=0$, and when $p^{-}=1$ it is $\Pr[\Ncount{\hat c}^{-}<n]=0$
because $\Ncount{\hat c}^{-}$, binomial $(n,1)$, equals $n$ almost surely. An
upper term is $\Pr[\Ncount{\hat c}^{+}>n]=0$ when $k^{+}(\hat c)=n$, and when
$p^{+}=0$ it is $\Pr[\Ncount{\hat c}^{+}>0]=0$ because
$\Ncount{\hat c}^{+}$, binomial $(n,0)$, equals $0$ almost surely.

\emph{(iv).} The left inequalities of \eqref{eq:width} are \eqref{eq:ordering}
(Lemma~\ref{lem:guard}\,(i)):
$q^{+}(\hat c)=k^{+}(\hat c)/n\ge\lceil np^{+}\rceil/n\ge p^{+}$ and
$q^{-}(\hat c)=k^{-}(\hat c)/n\le\lfloor np^{-}\rfloor/n\le p^{-}$.

For the upper one, $\kappa^{+}(\hat c)$ is the quantile \eqref{eq:kappaplus}
at $p^{+}$ and $\Ncount{\hat c}^{+}$ is binomial $(n,p^{+})$ (the
\tref{side-wise reading} of Definition~\ref{def:endpoint-rule}). Set
$t\coloneqq\sqrt{np^{+}(1-p^{+})/\varepsilon^{+}(\hat c)}>0$ and
$k\coloneqq\lceil np^{+}+t\rceil$. Any integer $j>k$ satisfies $j>np^{+}+t$,
so
\begin{equation*}
  \Pr[\Ncount{\hat c}^{+}>k]\ \le\ \Pr[\Ncount{\hat c}^{+}-np^{+}>t]
  \ \le\ \Pr\big[|\Ncount{\hat c}^{+}-np^{+}|\ge t\big]
  \ \le\ \frac{np^{+}(1-p^{+})}{t^{2}}\ =\ \varepsilon^{+}(\hat c)
\end{equation*}
by Lemma~\ref{lem:cheb}. If $k\le n$ this exhibits an admissible index in
\eqref{eq:kappaplus}, so $\kappa^{+}(\hat c)\le k$ by minimality; if $k>n$
then $\kappa^{+}(\hat c)\le n<k$ anyway. In both cases
$\kappa^{+}(\hat c)\le\lceil np^{+}+t\rceil$. Since also
$\lceil np^{+}\rceil\le\lceil np^{+}+t\rceil$, the guard gives
$k^{+}(\hat c)\le\lceil np^{+}+t\rceil\le np^{+}+t+1$, hence
$q^{+}(\hat c)-p^{+}\le t/n+1/n$, which is the stated bound because
$t/n=\sqrt{p^{+}(1-p^{+})/(n\varepsilon^{+}(\hat c))}$.

Symmetrically, $\kappa^{-}(\hat c)$ is the quantile \eqref{eq:kappaminus} at
$p^{-}$ and $\Ncount{\hat c}^{-}$ is binomial $(n,p^{-})$. Set
$t'\coloneqq\sqrt{np^{-}(1-p^{-})/\varepsilon^{-}(\hat c)}>0$ and
$k'\coloneqq\max\{0,\lfloor np^{-}-t'\rfloor\}$. If $\lfloor np^{-}-t'\rfloor
\le0$ then $k'=0$ and $\Pr[\Ncount{\hat c}^{-}<0]=0\le\varepsilon^{-}(\hat
c)$. Otherwise $k'=\lfloor np^{-}-t'\rfloor\le np^{-}-t'$, so any integer
$j\le k'-1$ satisfies $j<np^{-}-t'$ and
$\Pr[\Ncount{\hat c}^{-}<k']\le\Pr[\Ncount{\hat c}^{-}-np^{-}<-t']\le
\Pr[|\Ncount{\hat c}^{-}-np^{-}|\ge t']\le\varepsilon^{-}(\hat c)$ by
Lemma~\ref{lem:cheb}. Either way $k'$ is admissible in
\eqref{eq:kappaminus}, so $\kappa^{-}(\hat c)\ge k'$. In the unclipped case
$k'=\lfloor np^{-}-t'\rfloor>np^{-}-t'-1$; in the clipped case
$\lfloor np^{-}-t'\rfloor\le0$ forces $np^{-}-t'<1$, hence $np^{-}-t'-1<0=k'$.
In both cases $\kappa^{-}(\hat c)\ge k'>np^{-}-t'-1$. As
$\lfloor np^{-}\rfloor>np^{-}-1\ge np^{-}-t'-1$, the guard gives
$k^{-}(\hat c)\ge np^{-}-t'-1$ and hence $p^{-}-q^{-}(\hat c)\le t'/n+1/n$.

For the consequence: $\acSet\subseteq\prod_{\hat c}[q^{-}(\hat
c),q^{+}(\hat c)]$ by \eqref{eq:accset}, so every $\mathbf{q}\in\acSet$ has
$p^{-}_{\hat c}-q(\hat c)\le p^{-}_{\hat c}-q^{-}(\hat c)$ and
$q(\hat c)-p^{+}_{\hat c}\le q^{+}(\hat c)-p^{+}_{\hat c}$, which
\eqref{eq:width} bounds by $O(n^{-1/2})$ uniformly
over the finite alphabet, both bounds of \eqref{eq:width} applying at every
$\hat c$ because $0<p^{\mp}_{\hat c}<1$ and $\varepsilon^{\pm}(\hat c)>0$
there, with fixed values; $\mathsf{H}\subseteq\acSet$ is
Theorem~\ref{thm:box-construction}\,(iii). For the degenerate envelope
$p^{-}_{\hat c}=p^{+}_{\hat c}=\honest(\hat c)$, so the same two bounds give
$\|\mathbf{q}-\honest\|_{\infty}\le\max_{\hat c}\max\{q^{+}(\hat
c)-\honest(\hat c),\,\honest(\hat c)-q^{-}(\hat c)\}=O(n^{-1/2})$ for every
$\mathbf{q}\in\acSet$, and $\honest\in\acSet$ by Theorem~\ref{thm:box-construction}\,(iii),
so the other half of the Hausdorff distance is $0$.

\emph{(v).} Direct computation with the degenerate envelope at $p=\tfrac12$,
so that $\Ncount{}\coloneqq\Ncount{\hat c_0}^{-}=\Ncount{\hat c_0}^{+}$ is
binomial $(n,\tfrac12)$, and the exact rule. At
$n=2$ the law of $\Ncount{}$ is $(\tfrac14,\tfrac12,\tfrac14)$ on
$\{0,1,2\}$: $\Pr[\Ncount{}<0]=0\le\tfrac25$ and
$\Pr[\Ncount{}<1]=\tfrac14\le\tfrac25$ but
$\Pr[\Ncount{}<2]=\tfrac34>\tfrac25$, so $\kappa^{-}=1$;
$\Pr[\Ncount{}>0]=\tfrac34>\tfrac25$ and $\Pr[\Ncount{}>1]=\tfrac14\le\tfrac25$,
so $\kappa^{+}=1$. The guard is inert ($\lfloor1\rfloor=\lceil1\rceil=1$), so
$[k^{-},k^{+}]=[1,1]$ and $q(\hat c_0)=\tfrac12$, which on a two-letter
alphabet forces $\mathbf{q}=(\tfrac12,\tfrac12)$. At $n=3$ the law is
$(\tfrac18,\tfrac38,\tfrac38,\tfrac18)$ on $\{0,1,2,3\}$:
$\Pr[\Ncount{}<1]=\tfrac18\le\tfrac25$ but
$\Pr[\Ncount{}<2]=\tfrac12>\tfrac25$, so $\kappa^{-}=1$;
$\Pr[\Ncount{}>1]=\tfrac12>\tfrac25$ and
$\Pr[\Ncount{}>2]=\tfrac18\le\tfrac25$, so $\kappa^{+}=2$. The guard is again
inert ($\lfloor\tfrac32\rfloor=1$, $\lceil\tfrac32\rceil=2$), so
$[k^{-},k^{+}]=[1,2]$ and $q(\hat c_0)\in[\tfrac13,\tfrac23]$. The two sets are
as displayed and the inclusion is strict.

\emph{(vi).} On $\Chatset=\{\hat c_0\}$ the only vector with non-negative
coordinates summing to $1$ is $\delta_{\hat c_0}$, so
$\simplex{\Chatset}=\{\delta_{\hat c_0}\}$. Condition \eqref{eq:envelope}
gives ${\honest}^{+}(\hat c_0)\le1\le{\honest}^{+}(\hat c_0)$, that is,
${\honest}^{+}(\hat c_0)=1$, and ${\honest}^{-}(\hat c_0)\le1$, so
$\delta_{\hat c_0}\in\mathsf{H}\subseteq\simplex{\Chatset}$ and
$\mathsf{H}=\{\delta_{\hat c_0}\}$; for the degenerate envelope also
${\honest}^{-}(\hat c_0)=1$, that is, $\honest=\delta_{\hat c_0}$. Since
$\acSet\subseteq\simplex{\Chatset}$ by \eqref{eq:accset} and
$\acSet\supseteq\mathsf{H}\neq\varnothing$ by Theorem~\ref{thm:box-construction}\,(iii),
$\acSet=\simplex{\Chatset}=\{\delta_{\hat c_0}\}$ whatever the parameters. For
$\varepsilon^{\star}$: if $\Chatcon=\varnothing$, \eqref{eq:boxbounds} gives
the vacuous bounds $[0,1]$ and $\varepsilon^{\star}$ is an empty sum, hence
$0$. If $\Chatcon=\{\hat c_0\}$, the upper side has
$p^{+}_{\hat c_0}=1$, so Definition~\ref{def:endpoint-rule}\,(c) gives
$R^{+}=n$, the guard gives $k^{+}(\hat c_0)=n$ and the upper term
$\Pr[\Ncount{\hat c_0}^{+}>n]$ is $0$. If moreover
${\honest}^{-}(\hat c_0)=1$, the same convention gives $R^{-}=n$,
$k^{-}(\hat c_0)=\min\{n,n\}=n$, the box bounds
$q^{-}(\hat c_0)=q^{+}(\hat c_0)=1$ independent of every parameter, and the
lower term $\Pr[\Ncount{\hat c_0}^{-}<n]$ is $0$ because
$\Ncount{\hat c_0}^{-}$, binomial $(n,1)$, equals $n$ almost surely; so
$\varepsilon^{\star}=0$. The two cases exhaust $\Chatcon\subseteq\Chatset$,
and neither used $n$, $\ecomp$, the allocation or anything of $R$ beyond
Definition~\ref{def:endpoint-rule}\,(c). When ${\honest}^{-}(\hat c_0)<1$ the
lower term is the tail of a binomial $(n,{\honest}^{-}(\hat c_0))$ below
$k^{-}(\hat c_0)$, which Lemma~\ref{lem:guard}\,(ii) bounds by
$\varepsilon^{-}(\hat c_0)$ and which need not vanish; the bound of
Theorem~\ref{thm:completeness-bound} still holds, the true failure probability being
$\Pr[\operatorname{freq}_{\Chat_1^n}\notin\{\delta_{\hat c_0}\}]=0$.
\end{proof}
